%------------------------------------------------------------------------------%
% Luis A. D'Afonseca
%
%------------------------------------------------------------------------------%
\documentclass[a4paper,12pt,fleqn]{article}

\usepackage{style_avaliacao}

% \ShowAnswers

\newcommand{\D}[2]{\frac{\partial {#1}}{\partial {#2}}}
\newcommand{\DS}[2]{\frac{\partial^2 {#1}}{\partial {#2}^2}}
\newcommand{\DM}[3]{\frac{\partial^2 {#1}}{\partial {#2} \partial {#3}}}

\newcommand{\vetor}[2]{\left(\begin{array}{c} #1 \\ #2 \end{array}\right)}

\newcommand{\veci}{\bm{i}}
\newcommand{\vecj}{\bm{j}}
\newcommand{\veck}{\bm{k}}

%------------------------------------------------------------------------------%
\begin{document}

\makeHeader{CFVVI}{Prova Substitutiva -- Turma 4}{11/09/2024}

\textbf{O exercício correspondente a prova que será substituída vale 50 pontos}

% 01
%------------------------------------------------------------------------------%
\exercise{25}
Calcule o limite
\(
  \lim_{(x, y) \to (0, 0)} \frac{x^2 + y^2}{x^2 - y^2}
\)
ou mostre que ele não existe.

\begin{answer}
  Restringindo a função à reta \( y = 0 \) temos
  \[
    \left(\lim_{(x, y) \to (0, 0)} \frac{x^2 + y^2}{x^2 - y^2}\right) \evalat{y=0}{}
    = \lim_{x\to 0} \frac{x^2}{x^2}
    = \lim_{x\to 0} 1
    = 1
  \]
  Porém, se restringirmos $f$ à reta \( x = 0 \), temos
  \[
    \left(\lim_{(x, y) \to (0, 0)} \frac{x^2 + y^2}{x^2 - y^2}\right) \evalat{x=0}{}
    = \lim_{y\to 0} \frac{y^2}{-y^2}
    = \lim_{y\to 0} -1
    = -1
  \]
  Os limites ao longo de diferentes caminhos são diferentes, logo o limite não existe.
  \clearpage
\end{answer}

% 02
%------------------------------------------------------------------------------%
\exercise{25}
Encontre e classifique os pontos críticos da função
\(
  f(x, y) = x^4 + y^4 - 4xy
\)

\begin{answer}
  Calculando as derivadas parciais de $f$
  \[
    f_x(x, y) = 4x^3 - 4y \qquad\qquad f_y(x, y) = 4y^3 - 4x
  \]

  Como as derivadas parciais são polinômios elas existem em todos os pontos
  do plano. Portanto, temos apenas os pontos críticos onde as derivadas se anulam.
  Assim precisamos resolver o sistema \(\nabla f = 0\)
  \[
    \begin{cases}
      f_x(x, y) = 4x^3 - 4y = 0 \\
      f_y(x, y) = 4y^3 - 4x = 0
    \end{cases}
  \]
  que pode ser reescrito como
  \[
    \begin{cases}
      y = x^3 \\
      x = y^3
    \end{cases}
  \]
  Substituindo \( y = x^3 \) em \( x = y^3 \), obtemos
  \[
    x = \left(x^3\right)^3 = x^9
    \qquad\Rightarrow\qquad
    x^9 - x = 0
    \qquad\Rightarrow\qquad
    x(x^8-1) = 0
  \]
  temos duas soluções \(x=0\) ou \(x=1\).

  Para $x=0$ temos \(y=0\), enquanto para $x=1$ temos \(y=1\),
  portanto temos os pontos
  \[
    (x_1, y_1) = (0, 0)
    \qquad\qquad
    (x_2, y_2) = (1, 1)
  \]

  Para classificar os pontos críticos, usamos o teste da segunda derivada.

  Calculando as derivadas de segundas de $f$
  \[
    f_{xx}(x, y) = 12x^2 \qquad\qquad
    f_{yy}(x, y) = 12y^2 \qquad\qquad
    f_{xy}(x, y) = -4
  \]

  O discriminante é
  \[
    D(x, y) = f_{xx}(x, y)f_{yy}(x, y) - \big(f_{xy}(x, y)\big)^2
    = 12x^2 \; 12y^2 - (-4)^2
    = 144 x^2 y^2 -16
  \]

  Considerando o ponto \((x_1, y_1) = (0, 0)\).
  O discriminante é
  \[
    D(0, 0)
    = 144 \times 0^2 \times 0^2 - 16
    = -16
  \]
  Como \(D(0, 0) = -12 <0\) o ponto \((0, 0)\) é um ponto de sela

  Considerando o ponto \((x_2, y_2) = (1, 1)\)
  \[
    D(0, 0)
    = 144 \times 1^2 \times 1^2 - 16
    = 144 - 16
    = 128
  \]
  Como \(D(1, 1) = 128 >0\) o ponto é um máximo ou mínimo local,
  como \( f_xx(1, 1) = 12\times 1^2 = 12> 0 \) o ponto é um mínimo local.
  \clearpage
\end{answer}

% 03
%------------------------------------------------------------------------------%
\exercise{25}
Encontre as dimensões da caixa retangular fechada com máximo volume
que pode ser inscrita na esfera unitária.

\begin{answer}
  Queremos maximizar a função
  \[
    f(x, y, z) = (2x)(2y)(2z) = 8xyz
  \]
  sugeita à
  \[
    g(x, y, z) = x^2 + y^2 + z^2 = 1
  \]
  Sabemos também que $x$, $y$ e $z$ precisam sere positivos para que o volume $f$
  seja positivo.

  Para aplicarmos multiplicadores de Lagrange precisamos das derivadas
  parciais de $f$ e $g$
  \[
    f_x(x, y, z) = 8yz \qquad\qquad
    f_y(x, y, z) = 8xz \qquad\qquad
    f_z(x, y, z) = 8xy
  \]
  \[
    g_x(x, y, z) = 2x \qquad\qquad
    g_y(x, y, z) = 2y \qquad\qquad
    g_z(x, y, z) = 2z
  \]
  Vamos agora resolver o sistema \(\nabla f = \lambda\nabla g\) e \(g=1\)
  \[
    \begin{cases}
      8yz = 2\lambda x \\
      8xz = 2\lambda y \\
      8xy = 2\lambda z \\
      x^2 + y^2 + z^2 = 1
    \end{cases}
  \]
  Simplificando
  \[
    \begin{cases}
      4yz = \lambda x \\
      4xz = \lambda y \\
      4xy = \lambda z \\
      x^2 + y^2 + z^2 = 1
    \end{cases}
  \]

  Isolando o $\lambda$ na primeira equação
  \[
    \lambda = \frac{4yz}{x}
  \]
  e substituindo na segunda
  \[
    4xz
    = \lambda y
    = \frac{4yz}{x} y
    = \frac{4y^2z}{x}
    \qquad\Rightarrow\qquad
    x^2 = y^2
  \]
  Substituindo na terceira
  \[
    4xy = \lambda z
    = \frac{4yz}{x}z
    = \frac{4yz^2}{x}
    \qquad\Rightarrow\qquad
    x^2 = z^2
  \]
  Assim a terceira equação se torna
  \begin{align*}
    x^2 + y^2 + z^2 & = 1                      \\[2mm]
    3x^2            & = 1                      \\[2mm]
    x^2             & = \frac{1}{3}            \\[2mm]
    x               & = \frac{1}{\sqrt{3}}
  \end{align*}
  Portanto
  \[
    x = y = z = \frac{1}{\sqrt{3}}
  \]
\end{answer}

%------------------------------------------------------------------------------%
\end{document}
%------------------------------------------------------------------------------%
